\underline{Lema 3.11.} Hay un isomorfismo
\[
 D^2(A|R,M)\simeq
 \operatorname{Coker}\bigl(
 H^1(K_I,M)\longrightarrow
 \operatorname{Hom}_A(H_1(K_I),M)\bigr).
\]

\underline{Demostración.} Tenemos
\begin{align*}
D^2(A|R,M)
 &=H^1\bigl(A|R,
   \underline H^1_{A|R}(\operatorname{Der}_R(-,M))\bigr)\\
 &=\ker\left(
   H^1(P_A^{(1)}|R,M)
   \mathrel{\substack{\longrightarrow\\[-.6ex]\longrightarrow\\[-.6ex]\longrightarrow}}
   H^1(P_A^{(2)}|R,M)\right),
\end{align*}
porque
\[
 [\underline H^1_{A|R}(\operatorname{Der}_R(-,M))](P_A^{(0)})
 =[\underline H^1_{A|R}(\operatorname{Der}_R(-,M))](P)=0.
\]

Sea exacta
\[
 0\longrightarrow K\longrightarrow F\xrightarrow{\gamma}I
 \longrightarrow0.
\]
Denotemos por $P[F]$ el álgebra simétrica sobre $P$ generada por $F$.
Entonces $P[F]$ y $P[F]\otimes_PP[F]$ son álgebras de polinomios sobre
$R$, y existen morfismos sobreyectivos
\begin{align*}
P[F]&\longrightarrow P_A^{(1)}=P\times_AP,
&X&\longmapsto(0,\gamma(X)),\\
P[F]\otimes_PP[F]&\longrightarrow
P_A^{(2)}=P\times_AP\times_AP,
&X\otimes1&\longmapsto(0,\gamma(X),0),\\[-.3ex]
&&1\otimes X&\longmapsto(0,0,\gamma(X)),
\end{align*}
para $X\in F$. Pongamos
\[
 J=\ker(P[F]\longrightarrow P_A^{(1)}),
 \qquad
 J'=\ker(P[F]\otimes_PP[F]\longrightarrow P_A^{(2)}).
\]
Si $F_\gamma$ denota el ideal generado por
$\{X-\gamma(X)\mid X\in F\}$, entonces
\begin{align*}
J&=(F)\cap(F_\gamma)=(K)+(F\cdot F_\gamma),\\
J'&=(F\otimes1+1\otimes F)
 \cap(F_\gamma\otimes1+1\otimes F)
 \cap(F\otimes1+1\otimes F_\gamma)\\
&=(K\otimes1+1\otimes K)
 +(F\cdot F_\gamma\otimes1+1\otimes F\cdot F_\gamma)
 +(F\otimes F).
\end{align*}

%%%%%%%%%%  scan idx 96  |  folio 85  %%%%%%%%%%

Los tres morfismos
$\pi_i:P_A^{(2)}\longrightarrow P_A^{(1)}$ admiten levantamientos
$\widehat\pi_i:P[F]\otimes_PP[F]\longrightarrow P[F]$ dados, para
$X\in F$, por
\[
\begin{array}{c@{\qquad}c@{\qquad}c}
\widehat\pi_0:
\begin{cases}
 X\otimes1\longmapsto\gamma(X)-X,\\
 1\otimes X\longmapsto X,
\end{cases}
&
\widehat\pi_1:
\begin{cases}
 X\otimes1\longmapsto0,\\
 1\otimes X\longmapsto X,
\end{cases}
&
\widehat\pi_2:
\begin{cases}
 X\otimes1\longmapsto X,\\
 1\otimes X\longmapsto0.
\end{cases}
\end{array}
\]
Un cálculo directo muestra que
$\widehat\pi_0-\widehat\pi_1+\widehat\pi_2$ anula
\[
 K\otimes1+1\otimes K+1\otimes F\cdot F_\gamma
\]
y envía
$F\cdot F_\gamma\otimes1+F\otimes F$ sobre
$F\cdot F_\gamma$. Para todo
\[
 D\in\operatorname{Der}_R(P[F]\otimes_PP[F],M)
\]
se tiene
$D(F\cdot F_\gamma\otimes1+F\otimes F)=0$, porque $I\cdot M=0$.
En consecuencia,
\begin{align*}
D^2(A|R,M)
 &=\ker\left(
   D^1(P_A^{(1)}|R,M)
   \mathrel{\substack{\longrightarrow\\[-.6ex]\longrightarrow\\[-.6ex]\longrightarrow}}
   D^1(P_A^{(2)}|R,M)\right)\\
 &=\{[f]\in D^1(P_A^{(1)}|R,M)
       \mid f(F\cdot F_\gamma)=0\}\\
 &=\operatorname{Coker}\left(
   \operatorname{Der}_R(P[F],M)\longrightarrow
   \operatorname{Hom}_{P[F]}
   \left(\frac{(F)\cap(F_\gamma)}{F\cdot F_\gamma},M\right)
   \right)\\
 &=\operatorname{Coker}\left(
   \operatorname{Der}_P(P[F],M)\longrightarrow
   \operatorname{Hom}_{P[F]}
   \left(\frac{(F)\cap(F_\gamma)}{(F)(F_\gamma)},M\right)
   \right).
\end{align*}

Las sucesiones
\[
 0\longrightarrow(F)\longrightarrow P[F]
 \sisi{\overset{P[0]}{\relbar\joinrel\relbar\joinrel\relbar}}{\xrightarrow{P[0]}}P\longrightarrow0,
 \qquad
 0\longrightarrow(F_\gamma)\longrightarrow P[F]
 \sisi{\overset{P[\gamma]}{\relbar\joinrel\relbar\joinrel\relbar}}{\xrightarrow{P[\gamma]}}P_\gamma\longrightarrow0
\]
son exactas. Por tanto,
\[
 \frac{(F)\cap(F_\gamma)}{(F)(F_\gamma)}
 =\operatorname{Tor}_1^{P[F]}(P,P_\gamma)
 =H_1(K_I),
\]
donde $K_I$ es el complejo de Koszul asociado a
$F\xrightarrow{\gamma}I\subset P$. Obtenemos, pues,
\begin{align*}
D^2(A|R,M)
 &=\operatorname{Coker}\bigl(
   \operatorname{Hom}_P(F,M)\longrightarrow
   \operatorname{Hom}_P(H_1(K_I),M)\bigr)\\
 &=\operatorname{Coker}\bigl(
   H^1(K_I,M)\longrightarrow
   \operatorname{Hom}_A(H_1(K_I),M)\bigr).
\end{align*}
\hfill$\square$

Podemos enunciar ahora los hechos básicos que necesitaremos más
adelante, resumiendo los resultados ya establecidos.
